\chapter{Conclusion}
\label{ch:concl}

\section{Summary}
\label{sec:summary}

ViKi converts passive two-view RGB-D recordings of a human cube transfer
into separate hand, object and proposed UR10 trajectory artifacts. A
geometric plan is not itself a robot demonstration; dynamic replay with
free objects supplies the stricter simulation test. This thesis reports
the measured chain on 78 recordings and a local layout-conditioned
learning experiment rather than promising a later unmeasured result.

\section{Answers to the Research Questions}
\label{sec:answers}

\textbf{RQ1: pipeline yield.} Seventy-three of 78 recordings pass the
articulated hand gate; 44 of those pass the original object-quality gate
and 70 pass after independent scratch track recovery. Of 26 newly
recovered scenes, 22 pass geometric retargeting. Seventeen of those 22
also produce released stable stacks in the provisional MuJoCo replay.
The branch-specific denominators and the red-support contact in all 24
replayed scratch plans preclude one unqualified success percentage.
The external HO-Cap pilot measured 10.50--11.67~mm smoothed hand-joint MPJPE
against semi-automatic annotations on two short sequences using either the
selected pair or all eight views; it does not ground-truth the local cube
recordings or their object and robot trajectories. Its source episodes and
the local cohort were lost after the pilot and must be reacquired for a
reproducible final evaluation.

\textbf{RQ2: local demonstration generation.} One recovered human
motion generated 46 successful simulator variants in a 48-layout
$\pm20$~mm sweep and 27 in an independent 48-layout $\pm50$~mm sweep.
All derive from that one motion and touch the red support. This is
local, physics-filtered candidate generation, not arbitrary scene or
robot-configuration transfer.

\textbf{RQ3: learned correction.} A phase- and initial-layout-conditioned
model trained on 73 generated variants outperformed unchanged source
commands on the same 14 held-out teacher-selected layouts (14/14 versus
4/14) and 30 separately sampled layouts (30/30 versus 0/30). The earlier
plain state BC failed in all 17 closed-loop trials. The successful model
retains the source trajectory and observes privileged initial cube
positions; it is not RL or a vision-based closed-loop policy.

\section{Claim Boundary and Next Validation}
\label{sec:contrib-final}

The evidence supports a camera-to-simulation proof of concept for one
human source trajectory, local cube-layout changes, one old-UR10/Robotiq
model and one provisional contact configuration. It does not demonstrate
collision-free grasping, independent 3-D orientation accuracy,
source-disjoint generalization, visual-policy performance, benefit over
a separately trained no-data learner or real-robot transfer. Those
claims require additional tests rather than a larger number of
synthetic variants of the same source motion.

A small measurement-only set with known geometry, marked orientation
and independently held-out human trials is the practical next step
for the thesis evidence. Support-object clearance and contact-model
sensitivity need correction before treating a simulated success as a
robot-ready reference. A physical UR10 evaluation would additionally
require camera-to-robot registration and hardware safety validation;
it is not part of the present result.

\section{Concluding Remarks}
\label{sec:remarks}

The result is useful precisely because the stages are not conflated:
recording completeness, reconstruction coverage, model acceptance,
geometric IK, dynamic stack success and learning on new layouts are
reported separately. The local learning result is positive; the
red-support collisions and failed plain BC are equally important
boundaries on what ViKi has actually established.
